%%%%%%%%%%%%%%%%%%%%%%%%%%%%%%%%%%%%%%%%%
% Cleese Assignment (For Students)
% LaTeX Template
% Version 2.0 (27/5/2018)
%
% This template originates from:
% http://www.LaTeXTemplates.com
%
% Author:
% Vel (vel@LaTeXTemplates.com)
%
% License:
% CC BY-NC-SA 3.0 (http://creativecommons.org/licenses/by-nc-sa/3.0/)
% 
%%%%%%%%%%%%%%%%%%%%%%%%%%%%%%%%%%%%%%%%%

%----------------------------------------------------------------------------------------
%	PACKAGES AND OTHER DOCUMENT CONFIGURATIONS
%----------------------------------------------------------------------------------------

\documentclass[10pt]{article}

%%%%%%%%%%%%%%%%%%%%%%%%%%%%%%%%%%%%%%%%%
% Cleese Assignment
% Structure Specification File
% Version 1.0 (27/5/2018)
%
% This template originates from:
% http://www.LaTeXTemplates.com
%
% Author:
% Vel (vel@LaTeXTemplates.com)
%
% License:
% CC BY-NC-SA 3.0 (http://creativecommons.org/licenses/by-nc-sa/3.0/)
% 
%%%%%%%%%%%%%%%%%%%%%%%%%%%%%%%%%%%%%%%%%

%----------------------------------------------------------------------------------------
%	PACKAGES AND OTHER DOCUMENT CONFIGURATIONS
%----------------------------------------------------------------------------------------

\usepackage{lastpage} % Required to determine the last page number for the footer

\usepackage{graphicx} % Required to insert images

\setlength\parindent{0pt} % Removes all indentation from paragraphs

\usepackage[svgnames]{xcolor} % Enabling colors by their 'svgnames'
\usepackage[most]{tcolorbox} % Required for boxes that split across pages

\usepackage{booktabs} % Required for better horizontal rules in tables

\usepackage{listings} % Required for insertion of code

\usepackage{etoolbox} % Required for if statements

\usepackage{multido} % required for dotted lines

\usepackage[french]{babel} % English language hyphenation

\usepackage{svg}

%----------------------------------------------------------------------------------------
%	MARGINS
%----------------------------------------------------------------------------------------

\usepackage{geometry} % Required for adjusting page dimensions and margins

\geometry{
	paper=a4paper, % Change to letterpaper for US letter
	top=3cm, % Top margin
	bottom=3cm, % Bottom margin
	left=2.5cm, % Left margin
	right=2.5cm, % Right margin
	headheight=14pt, % Header height
	footskip=1.4cm, % Space from the bottom margin to the baseline of the footer
	headsep=1.2cm, % Space from the top margin to the baseline of the header
	%showframe, % Uncomment to show how the type block is set on the page
}

%----------------------------------------------------------------------------------------
%	FONT
%----------------------------------------------------------------------------------------
\usepackage{fontspec}
\usepackage{unicode-math}
\usepackage[utf8]{inputenc} % Required for inputting international characters
\usepackage[T1]{fontenc} % Output font encoding for international characters

\setmainfont[Ligatures=TeX]{Caladea}
%----------------------------------------------------------------------------------------
%	HEADERS AND FOOTERS
%----------------------------------------------------------------------------------------

\usepackage{fancyhdr} % Required for customising headers and footers

\pagestyle{fancy} % Enable custom headers and footers

% \lhead{\small\assignmentClass\ifdef{\assignmentClassInstructor}{\ (\assignmentClassInstructor):}{}\ \assignmentTitle} % Left header; output the instructor in brackets if one was set
% \chead{} % Centre header
% \rhead{\small\ifdef{\assignmentAuthorName}{\assignmentAuthorName}{\ifdef{\assignmentDueDate}{Due\ \assignmentDueDate}{}}} % Right header; output the author name if one was set, otherwise the due date if that was set

% \lfoot{} % Left footer
% \cfoot{\small Page\ \thepage\ sur\ \pageref{LastPage}} % Centre footer
% \rfoot{} % Right footer

% \renewcommand\headrulewidth{0.5pt} % Thickness of the header rule


% \fancypagestyle{plain}{
% 	\lhead{} % Left header; output the instructor in brackets if one was set
% 	\chead{} % Centre header
% 	\rhead{}

%   \renewcommand{\headrulewidth}{0pt}%
%   \fancyfoot[C]{\small Page\ \thepage\ sur\ \pageref{LastPage}}%
% }

%----------------------------------------------------------------------------------------
%	MODIFY SECTION STYLES
%----------------------------------------------------------------------------------------

% \usepackage{titlesec} % Required for modifying sections

%------------------------------------------------
% Section

% \titleformat
% {\section} % Section type being modified
% [block] % Shape type, can be: hang, block, display, runin, leftmargin, rightmargin, drop, wrap, frame
% {\Large\bfseries} % Format of the whole section
% {\assignmentQuestionName~\thesection} % Format of the section label
% {6pt} % Space between the title and label
% {} % Code before the label

% \titlespacing{\section}{0pt}{0.5\baselineskip}{0.5\baselineskip} % Spacing around section titles, the order is: left, before and after

%------------------------------------------------
% Subsection

% \titleformat
% {\subsection} % Section type being modified
% [block] % Shape type, can be: hang, block, display, runin, leftmargin, rightmargin, drop, wrap, frame
% {\itshape} % Format of the whole section
% {(\alph{subsection})} % Format of the section label
% {4pt} % Space between the title and label
% {} % Code before the label

% \titlespacing{\subsection}{0pt}{0.5\baselineskip}{0.5\baselineskip} % Spacing around section titles, the order is: left, before and after

% \renewcommand\thesubsection{(\alph{subsection})}

%----------------------------------------------------------------------------------------
%	CUSTOM QUESTION COMMANDS/ENVIRONMENTS
%----------------------------------------------------------------------------------------


%	TITLE SECTION
%----------------------------------------------------------------------------------------

% \newcommand{\authorstyle}[1]{{\large\usefont{OT1}{phv}{b}{n}#1}} % Authors style (Helvetica)

% \newcommand{\institution}[1]{{\footnotesize\usefont{OT1}{phv}{m}{sl}#1}} % Institutions style (Helvetica)

% \usepackage{titling} % Allows custom title configuration

% \newcommand{\HorRule}{\rule{\linewidth}{1pt}} % Defines the gold horizontal rule around the title

% \pretitle{
% 	\centering
% 	\vspace{-120pt} % Move the entire title section up
% 	\HorRule\vspace{10pt} % Horizontal rule before the title
% 	% \textbf
% 	\bfseries
% 	\fontsize{32}{36}
% 	\usefont{T1}{phv}{b}{n}
% 	\selectfont % Helvetica
% 	\color{DarkRed} % Text colour for the title and author(s)
% }

% \posttitle{\par\vskip 15pt} % Whitespace under the title

% \preauthor{} % Anything that will appear before \author is printed

% \postauthor{ % Anything that will appear after \author is printed
% 	% \vspace{10pt} % Space before the rule
% 	\par\HorRule % Horizontal rule after the title
% 	\vspace{-30pt} % Space after the title section
% }

 % Include the file specifying the document structure and custom commands

%----------------------------------------------------------------------------------------
%	ASSIGNMENT INFORMATION
%----------------------------------------------------------------------------------------

% Required
\newcommand{\assignmentQuestionName}{Question} % The word to be used as a prefix to question numbers; example alternatives: Problem, Exercise
\newcommand{\assignmentClass}{Physique Chimie} % Course/class
\newcommand{\assignmentTitle}{Activité n°1} % Assignment title or name
\newcommand{\assignmentAuthorName}{Chapitre 1} %

% Optional (comment lines to remove)
% \newcommand{\assignmentClassInstructor}{Jones 10:30am} % Intructor name/time/description
% \newcommand{\assignmentDueDate}{Monday,\ January\ 24,\ 2019} % Due date


%----------------------------------------------------------------------------------------
%	VARIABLES
%----------------------------------------------------------------------------------------

\newcommand{\titreActivite}{Activité 1 : Le Feu} % titre de l'activité
\newcommand{\objectif}{objectif de l'activité} % titre de l'activité
\newcommand{\contexte}{Lorem ipsum dolor sit amet, consectetur adipiscing elit. Suspendisse semper massa ac sem egestas, eget venenatis leo dapibus. Pellentesque habitant morbi tristique senectus et netus et malesuada fames ac turpis egestas. Sed vestibulum quam ex, id auctor leo feugiat ut. Pellentesque habitant morbi tristique senectus et netus et malesuada fames ac turpis egestas. Nunc nunc odio, tempor sit amet lorem non, consectetur dignissim justo. Donec ligula risus, finibus nec scelerisque ut, mattis mollis arcu. Donec eu neque nisl. Nulla facilisi.} % titre de l'activité
\newcommand{\resumeContexte}{Pourra-tu aider X et Y à faire Z?} % titre de l'activité




%----------------------------------------------------------------------------------------

\begin{document}
%----------------------------------------------------------------------------------------
%	TITLE PAGE
%----------------------------------------------------------------------------------------
\date{}
\title{\titreActivite}
\maketitle % Print the title page

% \thispagestyle{firststyle} % Suppress headers and footers on the title page


% \newpage

%----------------------------------------------------------------------------------------
%	QUESTION 1
%----------------------------------------------------------------------------------------

\underline{\textbf{Objectif}} :  \objectif

\vspace{10pt}

\underline{\textbf{Contexte}} :  \textit{\contexte}

\begin{center}
	\includegraphics[width=0.5\columnwidth]{activité.jpg} % Example image
\end{center}

\textbf{\resumeContexte}

%% ----------------------------------

\assignmentSection{Votre mission travail}

\begin{question}

\questiontext{Quelle est la vitesse de pointe d'une alouette du nord?}
\answerbox{4}

\end{question}

%----------------------------------------------------------------------------------------
%	QUESTION 2
%----------------------------------------------------------------------------------------

\begin{question}

\questiontext{À quel point une alouette de course courserait si elle pouvait courser?}

%--------------------------------------------

\begin{subquestion}{Supposons que ``courser" veut dire jeter.} % Subquestion within question

\answerbox{4}

\end{subquestion}

%--------------------------------------------

\begin{subquestion}{Supposons que ``courser" veut dire manger.} % Subquestion within question
\answerbox{4}
\end{subquestion}

%--------------------------------------------

\end{question}

%----------------------------------------------------------------------------------------
%	QUESTION 3
%----------------------------------------------------------------------------------------

\begin{question}

\questiontext{Identify the author of Equation \ref{eq:bayes} below and briefly describe it in English.}

\begin{equation}\label{eq:bayes}
	P(A|B) = \frac{P(B|A)P(A)}{P(B)}
\end{equation}

\answerbox{4}

\end{question}


%----------------------------------------------------------------------------------------

\assignmentSection{Pour aller plus loin}

%----------------------------------------------------------------------------------------
%	QUESTION 4
%----------------------------------------------------------------------------------------

\begin{question}

\questiontext{The table below shows the nutritional consistencies of two sausage types. Explain their relative differences given what you know about daily adult nutritional recommendations.
The table below shows the nutritional consistencies of two sausage types. Explain their relative differences given what you know about daily adult nutritional recommendations.
The table below shows the nutritional consistencies of two sausage types. Explain their relative differences given what you know about daily adult nutritional recommendations.
The table below shows the nutritional consistencies of two sausage types. Explain their relative differences given what you know about daily adult nutritional recommendations.
The table below shows the nutritional consistencies of two sausage types. Explain their relative differences given what you know about daily adult nutritional recommendations.
The table below shows the nutritional consistencies of two sausage types. Explain their relative differences given what you know about daily adult nutritional recommendations.
The table below shows the nutritional consistencies of two sausage types. Explain their relative differences given what you know about daily adult nutritional recommendations.}

\begin{table}[h]
	\centering % Centre the table
	\begin{tabular}{l l l}
		\toprule
		\textit{Per 50g} & Pork & Soy \\
		\midrule
		Energy & 760kJ & 538kJ\\
		Protein & 7.0g & 9.3g\\
		Carbohydrate & 0.0g & 4.9g\\
		Fat & 16.8g & 9.1g\\
		Sodium & 0.4g & 0.4g\\
		Fibre & 0.0g & 1.4g\\
		\bottomrule
	\end{tabular}
\end{table}

\answerbox{4}

\end{question}

\end{document}
